
\chapter[Implementation]{Implementation}{Implementation}\label{im}

\noindent
\rule{0.49\textwidth}{0.75pt} $_{\bigcirc}$ \rule{0.49\textwidth}{0.75pt}\\


\section{Implementation Tools, Data Handling Approach and System Limitations}
TrustChain implementation aims at data handling, with a focus on secure data handling as opposed to conventional data collection. React based forms are used to gather user inputs in the registration, guardian set-up and uploading of documents. The backend is done in Node.js and Express.js APIs. The upload of files is processed and sent to IPFS through Pinata that ensures decentralization. The Web3.js or Ethers.js is utilized to interact with blockchain and to enable the frontend to interact with the smart contracts to verify identity creation and access control. The frontend also takes care of the correct input validation prior to transmitting data to the backend services. Metamask, which uses wallet based authentication, enhances security and user ownership. It will not store sensitive data in the centralized servers. The API communication is designed such that it is consistent and reliable. This will provide ease of interaction among all layers of the architecture.

The system does not follow a conventional sample size approach as it is not data driven but identity driven. The platform does not deal with huge datasets but processes single identity records in real-time. Documents are uploaded by each user and unique hashed credentials are generated, used to verify the user again. Hashing mechanisms and decentralized storage are used to optimize data handling and reduce redundancy and enhance efficiency. The verification requests are dynamically executed in accordance with user permission and consent within smart contracts. The system has been planned to support multiple verification requests concurrently without any impact on the system. Processing with hash based processing provides the transfer of very little data between components. This save bandwidth and enhances efficiency in response. The verifications are independent to each other and this enhances scalability of operations. The architecture will be capable of extension in the future to accommodate more users.

Limitations:
The system implementation has certain constraints. The dependency on user supplied documents creates dependency on preliminary authenticity checks which are conceptually connected to government databases. Blockchain transaction and wallet interaction error handling can be complex to the non technical user. Also, the latency of networks in IPFS and blockchain activities can influence the response time. Large scale deployments can also be problematic due to the lack of rate limiting and advanced monitoring mechanisms and will have to be further optimized. In certain instances, wallet management may pose usability issues because users are dependent on it. Without appropriate recovery arrangements, loss of personal keys can result in difficulties in accessing them. It still remains in conceptual stage as it has to be integrated with other external systems such as government databases. Upgrades of smart contracts are tricky when they are deployed and need to be planned. These aspects underline the future areas of the improvement and development.

\section{Identity Data Structure and Cryptographic Components Description}
User Identity Data:

User related information like wallet address, decentralized identifier and linked guardian wallets. It also has fewer personal attributes such as a name, email and date of birth that are controlled to verify.
Applications: It is applied in authentication, ownership of identities and in controlling access rights in the system.

Credential Data:

Contains uploaded government identity documents on IPFS and the content identifiers generated. These credentials are never kept in the raw form on chain and are connected by hashed references.
Usage: It is applied to producing verifiable proofs, and retaining integrity of identity without revealing real documents.

Hash Data:

Represents multi layer hashed values of IPFS content identifiers using KECCAK 256 and secondary hashing. These hashes are secure representations of original documents.
Applications: The identity can be verified and compared during the process of verifying proof.



\section{Cryptographic Data Preparation and Processing}

The TrustChain system substitutes the traditional methods of preprocessing with cryptography processing method aimed at securing identity data prior to the verification process. After uploading identity documents, the system will directly turn them into secure representations in IPFS that creates a unique content identifier. This identifier gets hashed using KECCAK 256 and another hashing to make sure that it is irreversible. This stage of preparation is aimed at making sure that there is no raw data stored or relayed in the system. The resultant hashes are then formatted in such a way that they can be fed into Zero Knowledge Proof generation where ZK STARK algorithms are run over the hashed values to give verifiable proofs. With this method, there is no data cleaning or transformation like in the traditional systems. Rather, it is about providing the integrity of data, privacy and consistency of decentralized components. The ready cryptographic messages are then utilized throughout the verification lifecycle allowing the safe comparison, validation and access control without the disclosure of sensitive identity information.

\section{Description of Algorithms (Modular)}

Identity Check and Access Control:\\
Algorithm: ZK STARK (Zero Knowledge Scalable Transparent Argument of Knowledge) privacy preserving identity verification.\\
Usage: Checks user identity, creating cryptographic proofs of hashed credentials without disclosing any sensitive data. It assures the user that he or she is authentic and keeps all the data confidential throughout verification requests.\\

Hashing and Data Integrity:\\
Algorithm: KECCAK 256 multi layer with hashing.\\
Usage: Convert identifiers of IPFS generated content to secure irreversible hashes. These hashes serve as distinct identities of identity documents and are compared and verified without revealing raw data.\\

Decentralized Storage:\\
Algorithms: IPFS content addressing scheme.\\
Usage: Stores user uploaded identity-documents in a decentralized format and obtains a unique identifier of a content in a file. This guarantees high availability, fault tolerance and avoids reliance on centralized storage systems.\\

Smart Contract Verification:\\
Algorithm: Open source Ethereum Solidity based smart contract logic.\\
Usage: Runs the verification rules, handles the decentralized identifiers and implements the role based access control. It authenticates evidence and records of access requests, approvals and revocations that cannot be altered.\\

Wallet Based Authentication:\\
Algorithm: Public and Private Key Cryptography.\\
Usage: Signs wallets with Metamask or WalletConnect to authenticate users. It makes sure that the actions are initiated within the system by the rightful owner of the identity.\\

Access Control and Revocation:\\
Algorithm: Role Based Access Control based on smart contracts.\\
Usage: Controls permission states, period of access and revocation. It enables users to issue or revoke access on a dynamic basis with all the actions being logged transparently on the blockchain.\\

\section{Description of Tools for Implementation}

VS Code Visual Studio code:
Description: A source code editor created by Microsoft, which is light in weight but powerful.
Usage: Writing, debugging and administration of the frontend, back-end and smart contract code, offers an efficient development experience with blockchain and web technology extensions.

Git:
Description: A distributed version control system to track source code changes.
Usage: Controls versioning of an entire system during development, facilitating effective collaboration, tracking of code and retracting changes during various phases of the project.

GitHub:
Description: A web based platform for hosting and collaborating on Git repositories.
Usage: Manages the TrustChain project repository where people can collaborate, manage versions, track issues and maintain the project history with adequate documentation.

Postman:
Description: An API development and testing tool.
Usage: Tests and debugs API endpoints, ensuring they function correctly and meet the application's requirements.

MetaMask:
Description: A blockchain interaction and cryptocurrency wallet. Applications: Authenticate wallets, sign transactions and allow users to safely communicate with Ethereum smart contracts when verifying their identity and controlling access in an identity verification and access control operation.


\section{Interface Design}
The design of the platform interface of TrustChain is centered on simplicity, clarity and control by the user with secure interactions. The React.js frontend offers a simplistic dashboard to register, upload and access documents. Verification requests can be easily viewed, permission and track activity can be granted or revoked. An integration of wallets means that authentication is provided without involving complicated procedures. The system is designed along the human centered concepts and thus the system is reachable by both the technical and non technical users without compromising on efficiency and security during the flow of interactions.


\rule{0.49\textwidth}{0.75pt} $_{\bigcirc}$ \rule{0.49\textwidth}{0.75pt}\\
